\section{结论}
\label{sec:conclusion}

本文针对空地协同多功能雷达网络部署优化问题，提出了基于MOPSO-DT算法的扩展方法。主要贡献包括：

\begin{enumerate}
    \item 实现了基于Hertel-Mehlhorn的区域分解算法，有效处理复杂区域
    \item 采用垂直交点坐标变换，消除边界效应，边界ECR达到100\%
    \item 引入空地异构传播模型(A2G: $\alpha$=2.0, G2G: $\alpha$=4.0)，提高覆盖计算准确性
    \item 使用归一化目标函数增强Pareto解多样性，生成100个非支配解
    \item 揭示了覆盖率与干扰压制效能之间的强负相关关系(r=-0.912)
\end{enumerate}

在100km×100km区域、5部雷达的场景下，本文方法实现了100\%的ECR覆盖率和100\%的边界覆盖率，均超过论文基线的94.2\%和91.8\%。在更具挑战性的200km×200km、8部雷达场景下，算法生成了100个Pareto最优解，验证了方法的有效性。

\textbf{局限性：}更大规模场景(400km×400km)需要更多雷达才能达到理想覆盖率；Pareto解的多样性对目标函数缩放敏感。

\textbf{未来工作：}
\begin{itemize}
    \item 扩展到动态目标场景，考虑目标的移动性
    \item 引入更精细的地形遮挡模型，考虑非视距传播
    \item 研究多目标优化算法参数的自适应调整策略
\end{itemize}